\documentclass[10.5pt,a4paper]{article}
\usepackage{fontspec}
\usepackage[french]{babel}
\usepackage[margin=1.8cm,top=1.6cm]{geometry}
\usepackage{fontawesome5}
\usepackage{titlesec}
\usepackage{enumitem}
\usepackage[hidelinks]{hyperref}
\usepackage{tabularx}
\usepackage{needspace}

\hypersetup{
  pdftitle={CV — Mohamed Amine Charrada},
  pdfauthor={Mohamed Amine Charrada},
  pdflang={fr}
}

\pagestyle{empty}
\setlength{\parindent}{0pt}
\frenchsetup{StandardLists=true}

\titleformat{\section}{\large\bfseries}{}{0pt}{}[\vspace{-6pt}\rule{\textwidth}{0.6pt}]
\titlespacing*{\section}{0pt}{12pt}{8pt}

\newcommand{\entry}[2]{\noindent #1 \hfill #2\par}
\newcommand{\project}[1]{\needspace{5\baselineskip}\noindent\textbf{#1}\par\vspace{2pt}}
\setlist[itemize]{leftmargin=1.2em,itemsep=0pt,topsep=3pt,label=\textbullet}
\newenvironment{tasks}{\begin{itemize}[leftmargin=2.6em,topsep=3pt]}{\end{itemize}\vspace{3pt}}

\begin{document}

\begin{center}
  {\LARGE\bfseries MOHAMED AMINE CHARRADA}\\[4pt]
  Étudiant en cycle ingénieur informatique\\[3pt]
  \faPhone*\ +216 52 904 114 \quad
  \faEnvelope\ \href{mailto:mohamedamine.charrada@esprit.tn}{mohamedamine.charrada@esprit.tn} \quad
  \faMapMarker*\ Tunis, Tunisie\\
  \faLinkedin\ \href{https://www.linkedin.com/in/mohamed-amine-charrada-17a494315/}{LinkedIn : Mohamed Amine Charrada} \quad
  \faGithub\ \href{https://github.com/charrada-amine}{GitHub : charrada-amine}
\end{center}

\section*{Profil}
Élève ingénieur en informatique en 5\textsuperscript{e} année à ESPRIT, orienté vers le développement full stack
et le cloud. Expérience de bout en bout, de la conception du backend aux interfaces mobiles, web et desktop,
jusqu’au déploiement cloud et à l’automatisation CI/CD. Attaché à la qualité logicielle :
architecture modulaire, contrats d’API et tests automatisés.

\section*{Formation}
\entry{\textbf{ESPRIT — École Supérieure Privée d’Ingénierie et de Technologies}}{2022 – aujourd’hui}
Cycle ingénieur en informatique

\section*{Expériences professionnelles}
\entry{\textbf{Stagiaire en développement full stack — MEDCOM} (hybride)}{Juin 2026 – Oct. 2026}
\begin{tasks}
  \item Stage de 4 mois sur le projet Zennyt : application mobile, web et desktop.
  \item Conception et développement de mini-jeux d’évaluation, de bout en bout (Flutter et Spring Boot).
  \item Développement des API backend, des tests automatisés et de la console d’administration.
\end{tasks}
\entry{\textbf{Stagiaire en développement logiciel — ExpressDisplay} (sur site)}{Juin 2025 – Juil. 2025}
\begin{tasks}
  \item Développement de deux modèles de prédiction du temps en Python.
  \item Conception et développement d’une application de reporting (React Native).
  \item Développement d’une application mobile de gestion des ressources humaines.
  \item Participation à l’analyse des besoins et à la mise en place de solutions logicielles adaptées.
\end{tasks}
\entry{\textbf{Stage d’intégration en réseaux et cybersécurité — OACA} (sur site)}{Juil. 2023 – Août 2023}
\begin{tasks}
  \item Découverte de l’Office de l’Aviation Civile et des Aéroports et du fonctionnement de son service informatique.
  \item Étude de l’architecture réseau et des équipements d’infrastructure.
  \item Sensibilisation aux enjeux et aux bonnes pratiques de cybersécurité.
\end{tasks}

\section*{Projets}
\project{Zennyt — Plateforme de mise en relation entre candidats et recruteurs}
\begin{itemize}
  \item Application mobile, web et desktop : profils candidats, matching par swipe et Fit Score, messagerie.
  \item Mini-jeux d’évaluation des compétences, avec un score calculé côté serveur (anti-triche).
  \item Backend en monolithe modulaire et contrats d’API OpenAPI.
  \item Technologies : Flutter, Spring Boot (Java 21), PostgreSQL, Python (FastAPI), Docker.
  \item Hébergement sur Azure Container Apps.
\end{itemize}
\vspace{3pt}
\project{SamOps — Plateforme d’observabilité et d’optimisation cloud}
\begin{itemize}
  \item Conception et développement d’une plateforme DevOps orientée observabilité et optimisation des coûts cloud.
  \item Monitoring (métriques, logs), dashboards coûts/performance, intégrations cloud.
  \item Automatisation CI/CD et workers/agents pour analyser et améliorer la qualité du code.
  \item Architecture microservices (Go, NestJS, Python, etc.).
  \item Projet public déployé sur AWS (\href{https://samops.app}{samops.app}).
\end{itemize}
\vspace{3pt}
\project{Styleto — Application mobile Smart Fashion (iOS / Android)}
\begin{itemize}
  \item Application combinant IA et expériences en temps réel.
  \item Garde-robe numérique : détection, recommandations, essayage virtuel, marketplace, chat.
  \item Technologies : SwiftUI (iOS), Kotlin Jetpack Compose (Android), NestJS, MongoDB.
\end{itemize}
\vspace{3pt}
\project{Stage — Plateforme de recherche de stage assistée par des agents IA (PWA)}
\begin{itemize}
  \item Agent de veille qui collecte chaque jour les offres de stage et les structure automatiquement.
  \item Matching profil–offre avec re-classement par IA, génération de lettres et CV personnalisés.
  \item Coach d’entretien en mode vocal temps réel et assistant conversationnel outillé (agents IA).
  \item Extension Chrome de capture automatique des offres LinkedIn.
  \item Technologies : Next.js 16, Convex, Gemini (Vercel AI SDK), Better Auth, Polar.
\end{itemize}
\vspace{3pt}
\project{DinoBot — Harness d’orchestration d’agents IA de développement (macOS)}
\begin{itemize}
  \item Orchestrateur de 10 agents spécialisés (développeur, testeur QA, DevOps, critique…).
  \item Intégration des modèles : Claude Code, Codex et API LLM multi-fournisseurs.
  \item Runtime optimisé : sessions persistantes, cache de contexte, carte du projet, permissions et sandbox par agent.
  \item Technologies : Electron, React, Node.js, Swift.
\end{itemize}
\vspace{2pt}
\noindent\faGlobe\ Plus de projets sur mon portfolio : \href{https://www.charrada.me}{\textbf{charrada.me}}

\section*{Compétences}
\begin{tabular}{@{\hspace{0.6em}}l}
  \textbf{Langages :} Java, C, C++, Python, PHP, SQL, Go, Kotlin, Swift, JavaScript, TypeScript, Dart\\
  \textbf{Frameworks/SDK :} Spring Boot, NestJS, Next.js, Symfony, Flutter, React Native, FastAPI, Flask\\
  \textbf{Data \& UI :} NumPy, Pandas, React, JavaFX, Bootstrap\\
  \textbf{Bases de données/BaaS :} MySQL, PostgreSQL, MongoDB, Firebase, Supabase\\
  \textbf{DevOps/Cloud :} Git, Docker, AWS, Azure, GCP\\
  \textbf{Systèmes \& IoT :} Linux, Windows, Arduino, ESP32
\end{tabular}

\section*{Certifications}
\textbf{Hedera Certified Developer Associate (HCDA)} — Hedera\\
\textbf{Microsoft Applied Skills : Deploy cloud-native apps using Azure Container Apps} — Microsoft\\
\textbf{AWS Academy Graduate – Cloud Foundations} — Amazon Web Services

\section*{Leadership \& activités}
Membre — IEEE (2023)\\
Bénévole — Croissant-Rouge tunisien (depuis 2021)

\section*{Langues}
Arabe : langue maternelle \quad Français : B2 \quad Anglais : B2

\end{document}
